\documentclass[11pt,a4paper]{article}
\usepackage[margin=25mm]{geometry}
\usepackage[T1]{fontenc}
\usepackage{lmodern}
\usepackage{amsmath,amssymb,booktabs,graphicx,array,tabularx,xcolor}
\usepackage[hidelinks]{hyperref}
\setlength{\parindent}{0pt}
\setlength{\parskip}{5pt}
\setlength{\emergencystretch}{2em}
\renewcommand{\arraystretch}{1.12}
\newcommand{\code}[1]{\texttt{\detokenize{#1}}}
\newcommand{\fig}[3]{\begin{figure}[htbp]\centering
\includegraphics[width=\linewidth]{report_assets/#1}
\caption{#2}\label{#3}\end{figure}}
\newcommand{\tbl}[3]{\begin{table}[htbp]\centering\small
\caption{#2}\label{#3}#1\end{table}}
\makeatletter
\newcommand{\inlinecaption}[2]{\begingroup\def\@captype{#1}\caption{#2}\endgroup}
\makeatother
\newcommand{\MainTestTable}{%
\begin{tabular}{lrrrrrr}\toprule
Model & Total (\%) & CAGR (\%) & Arith. (\%) & Vol. (\%) & Sharpe & Max DD (\%)\\\midrule
M1 & 23.72 & 11.32 & 11.84 & 14.93 & 0.793 & -18.73\\
M2 & 43.25 & 19.86 & 20.88 & 23.52 & 0.888 & -26.15\\
M3 & 41.80 & 19.24 & 20.24 & 22.98 & 0.881 & -24.42\\
M4 & 50.83 & 23.01 & 22.91 & 20.96 & 1.093 & -19.04\\
SPY & 49.53 & 22.48 & 21.72 & 16.96 & 1.280 & -19.77\\
\bottomrule\end{tabular}}

\newcommand{\ValidationTable}{%
\begin{tabular}{lrrrrrr}\toprule
Model & CAGR 2021 & CAGR 2022 & CAGR 2023 & SR 2021 & SR 2022 & SR 2023\\\midrule
M1 & 29.18 & -13.39 & 9.45 & 1.855 & -0.408 & 0.571\\
M2 & 11.71 & -14.51 & 30.08 & 0.592 & -0.296 & 1.463\\
M3 & 12.37 & -11.19 & 24.89 & 0.594 & -0.168 & 1.227\\
M4 & 13.74 & -4.59 & 25.05 & 0.666 & 0.025 & 1.305\\
SPY & 31.09 & -19.57 & 26.59 & 2.290 & -0.780 & 1.815\\
\bottomrule\end{tabular}}

\newcommand{\AnnualTable}{%
\begin{tabular}{lrrrr}\toprule
Model & 2024 return (\%) & 2024 SR & 2025 return (\%) & 2025 SR\\\midrule
M1 & 16.39 & 1.331 & 6.30 & 0.441\\
M2 & 24.31 & 1.321 & 15.24 & 0.648\\
M3 & 24.54 & 1.279 & 13.86 & 0.624\\
M4 & 22.23 & 1.274 & 23.39 & 0.993\\
SPY & 26.94 & 1.965 & 17.79 & 0.910\\
\bottomrule\end{tabular}}

\newcommand{\PredictionTable}{%
\begin{tabular}{lrrrrr}\toprule
Model & IC & Rank IC & ICIR & $R^2$ & Coverage (\%)\\\midrule
M1 & 0.00206 & 0.00385 & 0.015 & -0.000092 & 99.9956\\
M2 & 0.01119 & 0.00813 & 0.074 & 0.000056 & 99.9956\\
M3 & 0.01074 & 0.00951 & 0.069 & 0.000022 & 99.9956\\
M4 & 0.01223 & 0.01195 & 0.092 & -0.000012 & 99.9956\\
\bottomrule\end{tabular}}

\newcommand{\CostTable}{%
\begin{tabular}{lrrrr}\toprule
Model & 0 bps CAGR & 5 bps CAGR & 10 bps CAGR & 0--5 drag (pp)\\\midrule
M1 & 14.38 & 11.32 & 8.34 & 3.06\\
M2 & 23.19 & 19.86 & 16.62 & 3.33\\
M3 & 22.56 & 19.24 & 16.02 & 3.32\\
M4 & 26.42 & 23.01 & 19.70 & 3.41\\
SPY & 22.54 & 22.48 & 22.42 & 0.06\\
\bottomrule\end{tabular}}

\newcommand{\ExecutionTable}{%
\begin{tabular}{lrrrrr}\toprule
Model & Turnover (\%) & Fees (USD) & Missing buys & Terminal events & Max trade/ADV (\%)\\\midrule
M1 & 10.76 & 61,603 & 1 & 3 & 0.119\\
M2 & 10.88 & 67,637 & 0 & 0 & 0.083\\
M3 & 10.89 & 67,356 & 0 & 0 & 0.085\\
M4 & 10.85 & 67,929 & 1 & 1 & 0.156\\
SPY & 0.20 & 1,248 & 0 & 0 & 0.003\\
\bottomrule\end{tabular}}

\newcommand{\SearchTable}{%
\begin{tabular}{lrrrrr}\toprule
Family & $\alpha$ & $R^2_{2021}$ & $R^2_{2022}$ & $R^2_{2023}$ & Mean $R^2$\\\midrule
Ridge & 100 & -0.000561 & -0.000284 & -0.000530 & -0.000458\\
Ridge & 1000 & -0.000463 & -0.000165 & -0.000207 & -0.000278\\
Ridge & 10000 & -0.000327 & -0.000108 & 0.000089 & -0.000115\\
Lasso & 0.0001 & -0.000453 & -0.000234 & -0.000370 & -0.000353\\
Lasso & 0.001 & -0.000107 & -0.000062 & 0.000121 & -0.000016\\
Lasso & 0.01 & 0.000174 & 0.000093 & 0.000127 & 0.000131\\
ElasticNet & 0.0001 & -0.000512 & -0.000271 & -0.000595 & -0.000459\\
ElasticNet & 0.001 & -0.000222 & -0.000123 & 0.000041 & -0.000101\\
ElasticNet & 0.01 & 0.000172 & -0.000030 & 0.000177 & 0.000106\\
\bottomrule\end{tabular}}

\newcommand{\AblationTable}{%
\begin{tabular}{lrrrrrr}\toprule
Estimator & $\alpha$ & $k$ & CAGR (\%) & Vol. (\%) & Sharpe & Max DD (\%)\\\midrule
Ridge (fixed) & 1000 & 48 & 23.01 & 20.96 & 1.093 & -19.04\\
Ridge (selected) & 10000 & 48 & 23.66 & 21.32 & 1.103 & -20.97\\
Lasso (selected) & 0.01 & 2 & 13.35 & 22.25 & 0.675 & -24.46\\
ElasticNet (selected) & 0.01 & 6 & 15.08 & 21.75 & 0.755 & -23.51\\
\bottomrule\end{tabular}}

\newcommand{\AblationPredictionTable}{%
\begin{tabular}{lrrrr}\toprule
Estimator & IC & Rank IC & $R^2$ & Turnover (\%)\\\midrule
Ridge & 0.01262 & 0.01151 & 0.000031 & 10.85\\
Lasso & 0.00592 & 0.00026 & 0.000040 & 10.87\\
ElasticNet & 0.00565 & 0.00119 & 0.000015 & 10.86\\
\bottomrule\end{tabular}}

\newcommand{\SparseTable}{%
\begin{tabular}{llr}\toprule
Estimator & Retained predictor & Coefficient\\\midrule
Lasso & reversal\_5d & 0.002024\\
Lasso & day\_return\_5 & -0.006793\\
ElasticNet & reversal\_5d & 0.004864\\
ElasticNet & day\_return\_5 & -0.008270\\
ElasticNet & relative\_dollar\_volume & -0.000494\\
ElasticNet & intraday\_efficiency & 0.002032\\
ElasticNet & intraday\_vwap\_reversal & 0.001796\\
ElasticNet & intraday\_rth\_return & -0.001690\\
\bottomrule\end{tabular}}

\newcommand{\TrainingTable}{%
\begin{tabular}{lrrrr}\toprule
Evaluation & Cutoff & Rows & Signal dates & Last training signal\\\midrule
validation\_2021 & 2020-12-31 & 379,033 & 754 & 2020-12-29\\
validation\_2022 & 2021-12-31 & 506,124 & 1006 & 2021-12-29\\
validation\_2023 & 2022-12-31 & 632,456 & 1257 & 2022-12-28\\
test\_2024\_2025 & 2023-12-29 & 757,943 & 1507 & 2023-12-27\\
\bottomrule\end{tabular}}

\newcommand{\VersionTable}{%
\begin{tabular}{ll}\toprule
Package & Version\\\midrule
numpy & 2.2.6\\
pandas & 2.3.3\\
pyarrow & 25.0.1\\
scikit-learn & 1.7.2\\
joblib & 1.6.0\\
threadpoolctl & 3.7.0\\
matplotlib & 3.10.9\\
\bottomrule\end{tabular}}

\newcommand{\RunTable}{%
\begin{tabular}{lp{0.77\linewidth}}\toprule
Configuration & Frozen result directory\\\midrule
M1 & \path{output/model1_20261107_223125_445454}\\
M2 & \path{output/model2_20261108_121514_275909}\\
M3 & \path{output/model3_20261108_134524_579273}\\
M4 & \path{output/model4_20261107_225220_083066}\\
Ablation & \path{output/model4_ablation_20261108_121514_275909}\\
\bottomrule\end{tabular}}


\begin{document}
\begin{titlepage}
\centering
{\large DSA5205 Project 2\par}
\vspace{16mm}
{\LARGE\bfseries Daily Stock Selection with Price--Volume Signals\par}
\vspace{5mm}
{\Large Intraday Information, Feature Enrichment,\par
and Regularization\par}
\vspace{12mm}
{\large Group 6\par}
\vspace{3mm}
{\normalsize
Wan Bicheng - A0356727R \\
Guo Jinghang - A0353941B \\
Zhao Liyan \\
Chen Yifan \\
\par}
\vspace{5mm}
{\normalsize 8 November 2026\par}
\vspace{12mm}
\begin{minipage}{0.92\linewidth}
\begin{abstract}
This report studies whether richer price--volume information improves daily
stock selection after transaction costs. Four Ridge configurations share the
same eligible stock universe, next-open return target, chronological evaluation,
and Top50Drop5 portfolio rule. The configurations use 18 intraday-derived
predictors, 16 basic daily indicators, 28 enriched daily factors, and the complete
48-factor set. The 28-factor configuration is an exact subset of the complete
model, allowing a nested comparison of the additional intraday-dependent
feature block. A separate nine-configuration ablation compares Ridge, Lasso,
and Elastic Net on the same 48 inputs using three expanding validation folds.
At 5 basis points per side, the complete fixed-penalty Ridge model earns a
2024--2025 CAGR of 23.01\%, compared with 19.86\% for the basic daily baseline,
19.24\% for the enriched daily model, 11.32\% for the intraday-only model, and
22.48\% for SPY. Its Sharpe ratio of 1.093 remains below SPY's 1.280.
Validation selects Lasso by prediction fit, but selected Ridge produces better
test-period portfolio performance. The results emphasize complementarity of
signals, the limits of feature expansion, and the distinction between
statistical accuracy and economic value.
\end{abstract}
\end{minipage}
\end{titlepage}

\section{Introduction and Research Questions}
\label{sec:intro}
This project compares four price--volume feature sets for daily stock
selection among historically eligible S\&P 500 constituents. All models use
Ridge regression, the same return target, and the same long-only portfolio
rule. SPY provides an investable benchmark, with US dollars as the accounting
currency. Intraday information is summarized from 5-minute bars, but every
strategy trades daily.

Daily returns describe momentum and reversal over several sessions. Intraday
returns and volume timing describe how the latest session unfolded. Combining
them may improve stock ranking, but additional predictors can also add noise.
Regularization helps control this problem by shrinking fitted coefficients.
Sparse forecasting motivates the Lasso comparison \cite{chinco}, while
research on complexity and shrinkage motivates retaining a broader set of
weak signals \cite{kelly}. Our experiment uses a smaller feature set and a
different trading horizon from those studies.

We ask three questions: Does enriching daily features improve a simple
baseline? Does adding intraday information improve the daily-only model?
On the same inputs, how do Ridge, Lasso, and Elastic Net compare? The main
comparison changes the features while fixing the learner; a separate
ablation changes the learner while fixing M4's 48 features.

\section{Market, Data Sources, and Cleaning}
\label{sec:data}
The source dataset covers historical S\&P~500 constituents from 2 January 2018
to 15 October 2026. Times use New York time and the NYSE calendar, with regular
sessions from 09:30 to 16:00 and early closes at 13:00. Prices are in US dollars.
Daily bars and 5-minute bars are downloaded separately: the models use the
prepared daily panel for daily predictors, and aggregate regular-session
5-minute bars for intraday predictors. The reported experiments end in 2025;
the later observations in the source dataset are not used in fitting,
validation, or testing.

\subsection{Original Providers and Lawful Data Access}
\label{sec:data-providers}
\begin{table}[htbp]\centering\small
\caption{Original data providers and acquisition channels.}
\label{tab:data-sources}
\begin{tabular}{p{0.20\linewidth}p{0.31\linewidth}p{0.39\linewidth}}
\toprule Input & Original provider & Acquisition channel\\\midrule
Index membership & S\&P Dow Jones Indices constituent changes &
Public \code{fja05680/sp500} repository, commit \code{a2430f2}\\
Daily and 5-minute bars & Exchange trades consolidated through the SIP &
Alpaca Market Data API v2, historical SIP feed\\
Adjustment factors & Vendor-applied corporate actions &
Alpaca raw, split-adjusted, and fully adjusted daily bars\\\bottomrule
\end{tabular}
\end{table}
The group obtained the data through the public membership repository and
Alpaca's authenticated API using its own free-plan credentials. The repository
version is recorded in \code{sp500/universe/source.json}.
The source repository is publicly accessible\footnote{\url{https://github.com/fja05680/sp500}};
Alpaca's historical SIP access without a subscription requires the query end
to be at least 15 minutes old\footnote{\url{https://docs.alpaca.markets/us/docs/market-data-faq}}.
This historical access differs from the free live IEX feed. Data were accessed
through permitted interfaces under the applicable usage conditions; credentials
are not part of the report or submitted source code.

\subsection{Sample Coverage and Security Identifiers}
\label{sec:data-coverage}
The collection spans 2,209 trading sessions and includes 365 calendar days
of pre-inclusion history and 90 days after removal, where available. These
buffers support lookbacks and exits; they do not expand the strategy's
membership eligibility. The download completed on 16 October 2026 contains
114,213,922 5-minute rows and 1,281,969 raw daily rows in 695 ticker files
per frequency, including SPY. The merged daily panel contains 1,280,621 rows
and 669 security identifiers, including the benchmark. These are collection
counts, not counts of eligible model observations.

A security master resolves ticker changes and reuse. A new symbol beginning
on the next trading session is linked to the previous symbol when its opening
price is within 6\% of the preceding close. There are 31 recorded rename links,
including FB to META, and three duplicate-history links. Large discontinuities,
such as UTX to RTX, are kept separate. Six reused symbols (IR, ARNC, FOXA, FOX,
CZR, and DOC) are split by issuer history. Securities use the latest ticker,
with a starting-year suffix where needed. This continuity rule is a practical
matching heuristic rather than an official permanent identifier.

\subsection{Historical Membership, Adjustments, and Missing Observations}
\label{sec:data-cleaning}
\paragraph{Membership and terminal returns.}
The historical file provides 698 membership spells, including 195 that ended
during the sample. A dated \code{in_index} flag restricts selection to historical
members rather than today's constituents. For 71 securities marked as delisted,
the prepared panel assigns terminal-return proxies: $-100\%$ for SIVB, SBNY,
and FRC, and zero for the other 68 cases, mostly acquisitions or other exits.
These are settlement assumptions, not verified realized proceeds. They remain
a limitation even though departed securities are retained in the dataset.

\paragraph{Adjustments.}
Let $s_t$ be split-adjusted close divided by raw close and $a_t$ be fully
adjusted close divided by raw close. The prepared panel uses $a_tP_t$ for
adjusted prices and $V_t/s_t$ for adjusted volumes. Factors are backward
adjusted to 15 October 2026 and chained across linked ticker changes.
Raw columns are retained. Intraday path summaries in the model code use raw
regular-session bars; daily return and execution calculations use the
panel's adjusted prices. These two input conventions are kept distinct.

\paragraph{Coverage and missing values.}
About 87\% of downloaded bars are in regular hours. The collection notes
report regular-session data on at least 99\% of relevant membership dates for
694 of the 695 tickers. In the prepared panel, 93.7\% of constituent-days have
all expected bars: normally 78, or 42 on a half day. Six dates between May 2018
and March 2022 have especially sparse coverage. Observed and expected bar
counts are recorded rather than filling missing raw bars. The source panel
also leaves its daily return field undefined across gaps longer than ten
sessions; the models construct their own next-open target.

The actual intraday feature builder requires at least 95\% bar coverage and
valid session endpoints. Otherwise its 18 intraday values are set to missing;
the stock-day can remain eligible, and training-median imputation handles
those missing predictors. There is no blanket deletion of all incomplete days.

\paragraph{Model eligibility and unusual observations.}
All models require historical membership, raw close of at least USD 5,
positive volume and adjustment factor, consistent OHLC bounds, 20-session
average dollar volume of at least USD 10 million, and 60 consecutive observed
closes. An adjustment warning excludes the current and following nine
sessions. Non-finite predictor values become missing and use training-only
imputation. There is no general winsorization; the residual-reversal/activity
interaction clips relative dollar volume to $[0.2,5]$. These filters and the
shared target ensure consistent samples across the four models.

\section{Experimental Design and Feature Architecture}
\label{sec:features}
\subsection{A Controlled Sequence of Four Models}
Table~\ref{tab:architecture} distinguishes the feature comparison from the
estimator-selection ablation. The main comparison uses the same Ridge
implementation with $\alpha_R=1000$ throughout. The eligible sample, target,
preprocessing rules, temporal folds, starting capital, portfolio construction,
and cost scenarios are shared. Changing the stocks, execution convention,
or trading intensity at the same time as the features would make return
differences difficult to interpret.

\begin{table}[htbp]\centering\small
\caption{Main feature configurations and their experimental roles.}
\label{tab:architecture}
\begin{tabular}{lp{0.45\linewidth}rl}\toprule
ID & Predictors & Count & Learner\\\midrule
M1 & Intraday aggregates only & 18 & Ridge, 1000\\
M2 & Basic daily price--volume indicators & 16 & Ridge, 1000\\
M3 & Enriched daily factors; exact daily-only subset of M4 & 28 & Ridge, 1000\\
M4 & M3 plus 18 intraday aggregates and two tail interactions & 48 & Ridge, 1000\\\bottomrule
\end{tabular}
\end{table}

The clearest nested comparison is M3 versus M4. Every M3 predictor has exactly
the same definition and value in M4, and the extra block contains exactly
20 variables. In contrast, M2 and M3 differ in definitions as well as feature
count. Their comparison evaluates basic versus enriched daily designs, not a
pure addition of twelve variables. M1 provides a complementary intraday-only
perspective, but it is not the same as trading every five minutes.

\subsection{M2: Basic Daily Indicators}
M2 uses 16 indicators: returns over one, five, and twenty sessions;
longer momentum windows that skip recent observations; daytime and overnight
returns; rolling volatility; daily range and closing position; deviations
from moving averages; and relative volume and dollar activity. These provide
an interpretable baseline covering trend, reversal, risk, and trading activity.
No 5-minute data or explicit interaction is included.
Appendix~\ref{app:basic} lists the formulas and lookbacks.

\subsection{M3: Enriched Daily Factors}
M3 uses the exact 28 daily-only predictors in M4. It adds market-residual
reversals, illiquidity, price-pressure measures, six Alpha101-style formulas,
and daily interactions. Residual reversal removes the return associated with
a lagged stock--SPY beta. Risk scaling and interactions let reversal signals
vary with volatility, liquidity, and activity. The formulaic factors follow
the approach of Kakushadze \cite{alpha101}; their exact definitions are given
in the appendix and the shared feature code.

M2 and M3 differ in definitions as well as feature count. Their comparison
therefore tests two daily feature designs, rather than simply adding twelve
columns to an otherwise unchanged model.

\subsection{M1 and M4: Intraday Information}
M1 uses 18 regular-session aggregates: opening and closing returns, volume
timing, realized variation, downside variation, path efficiency, price
location, and price--volume relationships. These distinguish sessions with
similar daily endpoints but different trading paths. The signed-volume
measure uses bar-return signs as a proxy; it is not observed order flow.

M4 combines M3's 28 daily factors with all 18 intraday aggregates and two
tail-return interactions. These interactions scale closing reversal by
intraday risk or closing activity. The M3--M4 comparison thus measures a
20-feature block, not the isolated effect of sampling frequency. M1 tests
whether the 18 intraday aggregates are useful without the daily context.

\section{Prediction Target, Timing, and Preprocessing}
\label{sec:labels}
\subsection{Decision versus Execution}
Let $t$ denote the signal session and $O^{a}_{i,t}$ and $C^{a}_{i,t}$ denote
adjusted open and close. The ordinary one-step target is
\begin{equation}
r_{i,t}=\frac{O^{a}_{i,t+2}}{O^{a}_{i,t+1}}-1.
\label{eq:label}
\end{equation}
The features $x_{i,t}$ use observations available by the close of $t$. The
portfolio is formed from the resulting scores, filled at $t+1$ open, and marked
at the next open. Information using the completed close of $t$ is not traded
at that same closing price.

\begin{table}[htbp]\centering\small
\caption{Information and execution timeline for a normal observation.}
\label{tab:timeline}
\begin{tabular}{lp{0.70\linewidth}}\toprule
Time & Operation\\\midrule
Close of $t$ & Complete the daily and intraday features; calculate and rank scores.\\
Open of $t+1$ & Rebalance using previously determined target weights and available quotes.\\
Open of $t+2$ & Value the holding interval and complete the ordinary return label.\\
Training cutoff & Admit an observation only if its label exit date is no later than the cutoff.\\\bottomrule
\end{tabular}
\end{table}

For a security with a terminal return marker, the label uses the source's
settlement proxy. This convention is shared by all models; a future terminal
event is not used to avoid an earlier purchase.

\subsection{Cross-Sectionally Standardized Training Labels}
The regression target is standardized within each signal date:
\begin{equation}
y_{i,t}=\frac{r_{i,t}-\bar r_t}{s_t},
\label{eq:ztarget}
\end{equation}
where $\bar r_t$ and $s_t$ are the mean and sample standard deviation of valid
eligible labels on that date. This focuses fitting on relative cross-sectional
outcomes and reduces changes in target scale across market conditions.
The fitted value is consequently a ranking score in standardized units,
not a calibrated expected percentage return.

Standardization is applied to observed training outcomes, not to the feature
inputs at prediction time. To prevent labels crossing a training boundary,
an observation is admitted only when its exit date is within the training
window. The final admitted signal is 27 December 2023, with exit on
29 December 2023.

\subsection{Training-Only Preprocessing}
For each training fold, non-finite feature values are treated as missing.
Missing entries are replaced by training medians, then each feature is
centered and scaled using its training mean and standard deviation.
The fitted transformations are applied unchanged to validation or test data.
Rows with missing predictors remain in the eligible sample, allowing the
four configurations to be compared on common stock-days.

\section{Regularized Linear Models}
\label{sec:models}
\subsection{Why Ridge?}
The models use the score $\hat y=b+z^\top\beta$, where $z$ contains the
standardized inputs. Ridge shrinks coefficients without requiring most of
them to be zero, which suits overlapping price--volume signals. Holding
Ridge fixed across M1--M4 lets us compare information sets directly.
M4 includes nonlinear feature transformations, but its fitted score remains
linear in the supplied columns.

\subsection{Objectives and Penalty Normalization}
For $n$ training rows, the implemented objectives are
\begin{align}
(\hat b,\hat\beta)_R
&=\arg\min_{b,\beta}
\sum_{j=1}^{n}(y_j-b-z_j^\top\beta)^2+\alpha_R\|\beta\|_2^2,
\label{eq:ridge}\\
(\hat b,\hat\beta)_L
&=\arg\min_{b,\beta}
\frac{1}{2n}\sum_{j=1}^{n}(y_j-b-z_j^\top\beta)^2+\alpha_L\|\beta\|_1,
\label{eq:lasso}\\
(\hat b,\hat\beta)_E
&=\arg\min_{b,\beta}
\frac{1}{2n}\sum_{j=1}^{n}(y_j-b-z_j^\top\beta)^2
+\alpha_E\rho\|\beta\|_1
+\frac{\alpha_E(1-\rho)}{2}\|\beta\|_2^2.
\label{eq:elastic}
\end{align}
These conventions match the installed scikit-learn estimators
\cite{sklearnridge,sklearnlasso,sklearnelastic}. The intercept is not the
regularized coordinate. All main models fix $\alpha_R=1000$; the separate
ablation searches penalties with $\rho=0.5$ for Elastic Net.

The raw alpha values are not directly comparable: Ridge uses summed squared
loss, whereas Lasso and Elastic Net divide the loss by $2n$. Lasso can set
coefficients exactly to zero; Elastic Net combines this sparsity with
quadratic shrinkage. All models include an unpenalized intercept.

Ridge uses scikit-learn's automatic solver. Lasso and Elastic Net use cyclic
coordinate descent with tolerance $10^{-6}$ and at most 20,000 iterations.
The iteration limit is checked. Coefficients describe contributions conditional
on the other inputs, rather than causal effects.

\section{Chronological Training, Validation, and Model Selection}
\label{sec:split}
\subsection{Expanding Validation and a Fixed Test-Period Fit}
The temporal design uses three expanding validation folds, followed by a
final 2018--2023 fit for 2024--2025. Table~\ref{tab:training}
records the admitted training samples, including the label-maturity
restriction. The same counts hold for the four main configurations.
\tbl{\TrainingTable}{Training windows and mature-label samples from run metadata.}
{tab:training}

Each fold fits preprocessing and coefficients on past observations only.
The final model is trained on mature 2018--2023 labels and held fixed through
2024--2025; features and portfolio scores still update daily.
Test signals run from 2 January 2024 to 29 December 2025. The first execution
is on 3 January 2024 and final valuation is on 31 December 2025, giving
500 return intervals.

Each validation year starts with a fresh portfolio. The test period is one
two-year portfolio path; its annual summaries do not restart the strategy.
Chronological folds, training-only preprocessing, mature labels, and
next-open execution control information flow.

\subsection{The Small Estimator-Selection Ablation}
The main comparison fixes Ridge alpha rather than tuning it
separately for each feature set. Estimator selection is isolated in
\code{model4_ablation.py}, using exactly M4's 48 inputs. Ridge tests
$\{100,1000,10000\}$; Lasso and Elastic Net each test
$\{0.0001,0.001,0.01\}$. Elastic Net fixes its L1 share at $0.5$ rather than
adding a second search dimension. There are nine candidates and
27 candidate-fold fits.

For candidate $c$ and validation year $v$, let $R^2_{c,v}$ be the
standardized-target prediction fit. The selection score is
\begin{equation}
S_c=\frac{1}{3}\sum_{v\in\{2021,2022,2023\}}R^2_{c,v}.
\label{eq:selection}
\end{equation}
The largest score selects the parameter within each family. Exact ties prefer
larger alpha. An overall validation-selected family is also recorded; exact
family ties prefer fewer average active features, then the stated
Ridge/Lasso/Elastic Net order. The choice uses validation results,
not test-period CAGR or Sharpe.

\tbl{\SearchTable}{Complete validation search on the 48-factor input set.
Values are dimensionless prediction $R^2$, not return percentages.}{tab:grid}
All three families choose the largest tested alpha: 10,000 for Ridge and
0.01 for Lasso and Elastic Net. Lasso has the largest mean score,
approximately 0.000131, followed by Elastic Net at 0.000106 and Ridge at
$-0.000115$. The preferred values lie at grid boundaries, so the experiment
identifies winners only within this grid, rather than a global optimum.

Each selected family is refitted on mature 2018--2023 labels and evaluated
in 2024--2025. All three selected estimators are reported. Validation results
are used for selection; the test comparison measures their subsequent
prediction and trading performance.

\section{From Predictions to a Self-Financing Portfolio}
\label{sec:trading}
\subsection{Top50Drop5 Selection}
At a signal close, eligible stocks with finite scores are ranked from highest
to lowest. The portfolio targets approximately 50 equal-weight long positions.
Existing holdings that lose eligibility must exit; vacancies are filled from
the current ranks. Among otherwise retained holdings, at most five
discretionary replacements are made: a better-ranked outsider replaces the
worst retained name only if its score is strictly higher.

The portfolio may differ from the unconstrained top 50 because voluntary
replacements are limited. Mandatory exits and refills do not consume the
five replacement slots. Retained positions are also rebalanced toward equal
weights, so five replacements do not imply only five orders. Turnover is
measured from actual simulated trades.

\subsection{Costs and the Cash Budget}
Every segment begins with USD 1 million. Let $q_{i,t}$ denote adjusted units,
$p_{i,t}$ the execution mark, and $B_t$ cash. Portfolio value is
\begin{equation}
V_t=B_t+\sum_i q_{i,t}p_{i,t}.
\end{equation}
For absolute traded dollar notional $A_t$, the charge is
\begin{equation}
C_t=c A_t,\qquad
A_t=\sum_i|\Delta q_{i,t}p_{i,t}|,\qquad
c=\frac{c_{\mathrm{bps}}}{10^4}.
\label{eq:fees}
\end{equation}
The main scenario is 5 bps per side; 0 and 10 bps provide sensitivity.
Both buys and sells are charged, including initial entry and final liquidation.
The solver determines fees and post-fee target allocations together, preventing
a fully invested long-only portfolio from borrowing simply to pay its fees.

One-sided turnover is $A_t/(2V_t^{\mathrm{pre}})$. The factor of two is only a
reporting convention: the fee rate multiplies the full $A_t$, not half of it.
Confusing those quantities would halve the intended cost. Cash interest and
the risk-free rate used for Sharpe are zero. These long-only configurations have
no short borrowing charge.

\subsection{Missing Quotes and Settlement}
An unfilled purchase leaves its allocation in cash. Existing holdings without
a tradable quote are frozen and valued using available marks. A held security
with a terminal marker is settled using the source proxy. Remaining tradable
positions are liquidated at the final valuation, with fees charged.

The backtest uses fractional adjusted units and idealized opening fills.
It does not simulate order queues, partial fills, or market impact. The three
flat-cost scenarios therefore describe execution assumptions rather than
observed trading costs.

\section{Evaluation Measures and Verification}
\label{sec:metrics}
\subsection{Forecast Metrics}
On each date, Pearson IC correlates predicted scores with valid realized
returns, and Spearman Rank IC correlates their ranks. We report their daily
means and define ICIR as mean IC divided by its sample standard deviation,
without annualization. Constant score or target cross-sections do not receive
a fabricated correlation; they are treated as undefined.

The prediction-fit measure is
\begin{equation}
R^2_{\mathrm{zero}}=
1-\frac{\sum_{i,t}(y_{i,t}-\hat y_{i,t})^2}
{\sum_{i,t}y_{i,t}^2},
\label{eq:r2}
\end{equation}
where $y$ is the ex-post standardized label from Equation~\eqref{eq:ztarget}.
The benchmark is a zero standardized forecast, not a raw-return forecast
and not the return to holding SPY. Valid-label coverage measures how many
predicted eligible rows have an observable target. It is not a measure of
coverage of all listed stocks or completeness of every 5-minute session.

Ranking and squared-error fit measure different properties: the regression
loss covers all stock-days, whereas the portfolio buys a limited set of
high-ranked stocks and retains some earlier holdings.

\subsection{Net Return and Risk Measures}
For $T$ observed daily intervals with net returns $R_t$, the report uses
\begin{align}
\mathrm{TotalReturn}&=\frac{V_T}{V_0}-1,&
\mathrm{CAGR}&=\left(\frac{V_T}{V_0}\right)^{252/T}-1,\\
\mu_{\mathrm{ann}}&=252\bar R,&
\sigma_{\mathrm{ann}}&=\sqrt{252}\,s(R_t),\\
\mathrm{Sharpe}_{0}&=\frac{\mu_{\mathrm{ann}}}{\sigma_{\mathrm{ann}}},&
\mathrm{MDD}&=\min_t\left(\frac{V_t}{\max_{u\le t}V_u}-1\right).
\end{align}
The drawdown path includes initial capital. Volatility uses the sample
standard deviation. Drawdown is based on daily valuation marks, not the worst
intraday loss. CAGR and annualized arithmetic return are reported separately
because they answer different questions.

SPY uses the same execution/valuation dates and fee rate. It is
bought and held, so its turnover and fee burden are much lower than those of
the selection strategies. This is an economically relevant difference between
active trading and the investable alternative, not a reason to force SPY to
rebalance daily.

Saved portfolio returns are checked against final wealth and fees. Report
figures are reconstructed from saved predictions and checked against the
stored metrics. M3's 28 features also match the corresponding M4 columns,
including row order and labels.

\section{Results: Feature Sets under a Common Ridge Learner}
\label{sec:mainresults}
\subsection{Validation Stability}
\tbl{\ValidationTable}{Validation CAGR in percent and Sharpe ratio at 5 bps per side.
Each year is a separate forward fold.}{tab:validation}


Table~\ref{tab:validation} shows substantial
period dependence. In 2021, M1 earns 29.18\% CAGR and Sharpe 1.855, close to
SPY's 31.09\% CAGR but below its 2.290 Sharpe. M2--M4 have much lower returns
that year. Intraday-only features are therefore not uniformly weak: their
strongest validation year differs from their later behavior.

In 2022, all four models lose money in CAGR terms, but each loses less than
SPY's 19.57\%. M4 is the strongest active model at $-4.59$\%, compared with
$-11.19$\% for M3, $-13.39$\% for M1, and $-14.51$\% for M2. The protection is
relative, not absolute. The strategy is not market-neutral, hedged, or designed
to guarantee a positive return in a falling market.

In 2023, M2 leads the active configurations with 30.08\% CAGR. M3 and M4 earn
about 24.89\% and 25.05\%, while M1 earns 9.45\%. M4 does not dominate
every validation fold, and the enriched daily model does not always beat
the basic one. Averaging away these differences would conceal an important
part of the experiment.



\subsection{The 2024--2025 Test Period}
\tbl{\MainTestTable}{Test performance at 5 bps per side over 500 return intervals.
Arith.\ is annualized arithmetic return; Vol.\ is annualized volatility.
All return, volatility, and drawdown entries are percentages.}{tab:test}
\fig{main_performance.png}{Common-date test net wealth and drawdown for the four main models
and SPY, including entry and final liquidation costs.}{fig:test}

The main ranking by CAGR is M4, M2, M3, and M1. M4 earns 50.83\% total return
and 23.01\% CAGR, versus SPY's 49.53\% and 22.48\%. The annualized return
advantage over SPY is only about 0.54 percentage points. It is a modest
difference in the observed sample, not strong evidence of persistent alpha.

M4's Sharpe is 1.093, below SPY's 1.280, because annualized volatility
is 20.96\% versus 16.96\%. Its drawdown of $-19.04$\% is slightly shallower
than SPY's $-19.77$\%. Return and drawdown alone do not fully characterize
risk-adjusted performance: the model can have a slightly better worst
drawdown yet a worse ratio of average return to overall variability.

M1 earns 11.32\% CAGR with 14.93\% volatility, the lowest volatility of the
active models. Its drawdown of $-18.73$\% is also the shallowest of the four.
However, its Sharpe of 0.793 and markedly lower return do not establish an
attractive overall substitute for SPY. The intraday-only representation trades
off much of the observed return for a different risk path.

M2 and M3 produce similar performance. M2 has a 19.86\% CAGR and 0.888 Sharpe,
while M3 has 19.24\% and 0.881. M3 reduces volatility from 23.52\% to 22.98\%
and drawdown from $-26.15$\% to $-24.42$\%. Enrichment offers some risk
reduction in this sample but does not yield a clear improvement in the overall
return--risk profile.

\subsection{The Nested M3--M4 Comparison}
Adding the 20 intraday-dependent predictors to M3 raises CAGR from 19.24\%
to 23.01\%, about 3.77 percentage points. Sharpe rises by about 0.21,
annualized volatility falls by about 2.02 percentage points, and drawdown
magnitude falls by about 5.38 percentage points. These changes occur under
the same fixed Ridge alpha and portfolio rule.

This supports an incremental contribution of the combined intraday block
within this particular model. It does not show that high-frequency inputs
are useful in isolation: M1 performs considerably worse. Nor does it establish
whether improvement comes primarily from one intraday aggregate, the two
tail interactions, or their joint influence on selected stocks.
The reported comparison measures their combined contribution.



\subsection{Annual Test Outcomes}
\tbl{\AnnualTable}{Annual subdivisions of the test path at 5 bps per side.
Returns are realized segment totals; SR is the annualized Sharpe within that segment.}
{tab:annual}
In 2024, M2 and M3 return 24.31\% and 24.54\%, above M4's 22.23\%, while
SPY returns 26.94\%. In 2025, M4 returns 23.39\%, above M2's 15.24\%,
M3's 13.86\%, and SPY's 17.79\%. M4's full-period advantage is not a uniform
lead at every annual boundary.

M1's annual returns fall from 16.39\% in 2024 to 6.30\% in 2025, with Sharpe
falling from 1.331 to 0.441. Its strong 2021 validation performance is not a
stable property of the intraday feature design. The yearly split is descriptive
and based on exit dates. It is not a separate model-selection exercise or
a claim about future regimes.

\subsection{Forecast Fit versus Ranking}
\tbl{\PredictionTable}{Test-period forecast metrics for the main models.
ICIR is not annualized; $R^2$ is dimensionless.}{tab:prediction}


M4 has the highest test Rank IC, approximately 0.01195, compared with 0.00951
for M3, 0.00814 for M2, and 0.00385 for M1. Its Pearson IC and ICIR also exceed
the daily-control values. These differences are directionally consistent
with the portfolio ranking, although no claim of statistical significance is made.

Squared-error fit gives a different ordering. M2 and M3 have small positive
$R^2$, approximately 0.000056 and 0.000022, while M4's is slightly negative
at about $-0.000012$. M4 can generate a stronger constrained portfolio
despite failing to beat the zero standardized forecast in pooled squared
error. Whole-cross-section error magnitude and tail-sensitive portfolio
rankings are different objectives.

All four models have approximately 99.9956\% valid-label coverage among
predicted eligible rows, so their scores use a common evaluation sample.

\subsection{Turnover, Costs, and Execution Diagnostics}
\tbl{\CostTable}{Test CAGR in percent at three per-side cost scenarios.
Drag is the zero-cost minus 5-bps CAGR difference in percentage points.}{tab:costs}
\fig{costs.png}{Cost sensitivity for the main feature comparison and the separate
48-factor estimator ablation. Each point uses the actual cash ledger at that cost.}
{fig:costs}


Cost sensitivity is economically material even though the learner changes
only the ranking. M4's zero-cost CAGR is about 26.42\%, falling to 23.01\%
at 5 bps and 19.70\% at 10 bps. At 10 bps it no longer exceeds SPY's CAGR.
The apparent active advantage depends on the assumed execution environment
and is not robust to all flat-cost scenarios.

Active portfolios have similar average one-sided turnover, around
10.8--10.9\% per session, while SPY's is about 0.20\% after averaging entry
and exit over the test interval. M4 incurs approximately USD 67,929 in
simulated fees on the USD 1 million starting-capital path. The dollar total
is path dependent, not a fixed annual expense ratio. Differences between
zero- and positive-cost CAGR also reflect compounding and fee-dependent allocations.

All four models finish with zero unliquidated positions. M1 and M4 each have
one missed planned purchase; no frozen-holding days occur in the test runs.
The backtest does not enforce a volume-participation cap or model market
impact, so these results do not establish live trading capacity.

\section{Results: Regularization on the Same 48 Factors}
\label{sec:ablation}
\subsection{Selected Estimators and the Fixed Reference}
\tbl{\AblationTable}{Test results for the 48-factor regularization ablation at 5 bps.
$k$ counts coefficients with absolute value above $10^{-12}$.
The fixed reference is standalone M4, not the selected Ridge.}{tab:ablation}
\fig{ablation_performance.png}{Selected regularizers on identical 48-factor inputs,
with original fixed Ridge as a dashed reference and SPY as the investable benchmark.}
{fig:ablation}

Selected Ridge uses alpha 10,000 and earns 23.66\% CAGR with Sharpe 1.103.
Selected Lasso and Elastic Net use alpha 0.01 and earn 13.35\% and 15.08\%
CAGR, with Sharpes of 0.675 and 0.755. Their drawdowns, $-24.46$\% and
$-23.51$\%, are deeper than selected Ridge's $-20.97$\%.

Selected Ridge improves CAGR slightly over fixed M4, from 23.01\% to
23.66\%, but also has higher volatility and a deeper drawdown. Its Sharpe
gain is small. It is not an unqualified improvement on every risk dimension.
The main feature comparison remains fixed at alpha 1000 so these changes
are not mixed into M3 versus M4.

Lasso is the validation-selected family under Equation~\eqref{eq:selection}.
Validation prediction fit selected Lasso, while the test-period trading
comparison favored Ridge. All three selected estimators are reported to
compare prediction accuracy with portfolio performance.

\subsection{Sparsity and Retained Signals}
\tbl{\SparseTable}{Nonzero standardized coefficients in the final sparse estimators.
Other candidate coefficients are zero under the reporting threshold.}{tab:sparse}


Ridge retains all 48 coordinates. Lasso retains only \code{reversal_5d}
and \code{day_return_5}. Its final solution discards every intraday aggregate
even though the candidate set contains them. Elastic Net retains six variables:
those two daily signals, relative dollar volume, intraday efficiency,
intraday VWAP reversal, and intraday regular-session return.

Sparse coefficients do not imply low turnover. Each retained predictor varies
across stocks and dates, and the portfolio rule still rebalances positions.
The sparse strategies average roughly 10.9\% daily one-sided turnover,
similar to Ridge.

\subsection{Why Prediction-Based Selection and Trading Diverge}
\tbl{\AblationPredictionTable}{Test prediction and turnover diagnostics for
selected regularizers. $R^2$ is dimensionless; turnover is percent.}{tab:ablationprediction}
Lasso has test $R^2$ of approximately 0.000040, slightly above selected
Ridge's 0.000031. Yet its Rank IC is only about 0.00026, compared with
Ridge's 0.01151. Elastic Net's Rank IC is also much lower, about 0.00119.
Sparse estimators can reduce broad squared error through shrinkage while
providing less useful ordering for the upper-tail portfolio.

The prediction loss does not account for portfolio replacement limits,
transaction costs, or the concentration of mistakes among selected stocks.
This helps explain why a better prediction $R^2$ need not produce a better
portfolio. Selecting by validation CAGR or Sharpe would test a different
criterion from the one used here.

\section{Discussion and Limitations}
\label{sec:limits}
The clearest result is the improvement from M3 to M4 under a common learner
and trading rule. Intraday-only M1 performs much worse, suggesting that daily
context and intraday summaries are more useful together. The 20-feature block
also contains two interactions, so their effects are evaluated jointly with
the 18 intraday aggregates.

M3 does not clearly improve on the simpler M2 baseline, and sparse
regularization does not improve the selected portfolio. Adding features or
removing coefficients therefore offers no automatic performance gain.
These findings apply to the specified features, small parameter grid, and
portfolio rule.

Only three validation years and two test years are evaluated. No formal
significance test or confidence interval is reported for model differences.
M4's small CAGR advantage over SPY should be read alongside its lower Sharpe
and cost sensitivity. A return difference from SPY is not a risk-adjusted
alpha estimate.

Opening-price fills, fractional units, flat costs, and terminal settlement
proxies simplify execution. Drawdown uses daily marks and misses intraday
extremes. The reported results use fixed test-period coefficients and
prediction-based validation selection.

\noindent\begin{minipage}{\linewidth}
\section{Reproducibility}
\label{sec:repro}

Under the group's collection setup, Alpaca's API rate limits make downloading
S\&P~500 data take approximately 7--11 hours. Therefore, we provide the prepared S\&P~500 dataset for download.


Download the prepared \code{sp500} dataset using the shared download link\footnotemark.
Extract the archive into the project root, alongside \code{model1.py},
so that \code{sp500/panel/daily.parquet} and
\code{sp500/raw/alpaca/} are available at those relative paths.

From the project root, the main commands are:
\begin{verbatim}
python -m pip install -r requirements.txt
python model4.py
python model1.py
python model2.py
python model3.py
python model4_ablation.py
python utils/build_report_assets.py --latest
\end{verbatim}
\end{minipage}
\footnotetext{Prepared S\&P~500 dataset: \href{https://nusu-my.sharepoint.com/shared?listurl=https\%3A\%2F\%2Fnusu\%2Dmy\%2Esharepoint\%2Ecom\%2Fpersonal\%2Fliyanzhao\%5Fu\%5Fnus\%5Fedu\%2FDocuments\&id=\%2Fpersonal\%2Fliyanzhao\%5Fu\%5Fnus\%5Fedu\%2FDocuments\%2Fdata\%5Fshare\&shareLink=1\&ga=1}{\textcolor{blue}{SharePoint dataset download}}.}


The model scripts fit and evaluate each configuration. The asset builder
generates verified tables and figures from completed runs. Use \code{--rebuild}
on a model script to regenerate its features.

Each timestamped \code{output} folder contains models, metrics, predictions,
two figures, metadata, and notes. Features, fitting, and portfolio accounting
are implemented in \code{utils/features.py}, \code{utils/experiment.py}, and
\code{utils/backtest.py}, respectively.

\tbl{\VersionTable}{Recorded package versions. The declared seed is zero;
solvers use deterministic ordering.}{tab:versions}
Run identifiers, timestamps, and settings are recorded in each output folder's
metadata. The prepared-data download supports the commands above. Rebuilding
directly from providers additionally requires the original download and
security-master/panel-build scripts.

\section{Conclusion}
Combining daily and intraday-derived signals gives the strongest portfolio
among the four fixed-Ridge models. M4 earns 23.01\% test CAGR at 5 bps per
side, compared with 19.86\% for M2 and 19.24\% for M3. Its 1.093 Sharpe
remains below SPY's 1.280, and its modest CAGR advantage disappears at 10 bps.

Lasso wins the validation prediction comparison but selected Ridge produces
better test trading performance. The project therefore shows both the value
of complementary inputs and the need to evaluate forecasts through an
explicit portfolio and transaction-cost rule.

\clearpage
\begin{thebibliography}{9}
\bibitem{chinco}
Chinco, A., Clark-Joseph, A. D., and Ye, M. (2019).
Sparse Signals in the Cross-Section of Returns.
\emph{The Journal of Finance}, 74(1), 449--492.
\url{https://doi.org/10.1111/jofi.12733}.
\bibitem{kelly}
Kelly, B. T., Malamud, S., and Zhou, K. (2024).
The Virtue of Complexity in Return Prediction.
\emph{The Journal of Finance}, 79(1), 459--501.
\url{https://doi.org/10.1111/jofi.13298}.
\bibitem{alpha101}
Kakushadze, Z. (2016). 101 Formulaic Alphas.
\emph{Wilmott}, 2016(84), 72--81.
\url{https://arxiv.org/abs/1601.00991}.
\bibitem{sklearnridge}
Scikit-learn developers. Ridge, version 1.7.2 documentation.
\url{https://scikit-learn.org/1.7/modules/generated/sklearn.linear_model.Ridge.html}.
\bibitem{sklearnlasso}
Scikit-learn developers. Lasso, version 1.7.2 documentation.
\url{https://scikit-learn.org/1.7/modules/generated/sklearn.linear_model.Lasso.html}.
\bibitem{sklearnelastic}
Scikit-learn developers. ElasticNet, version 1.7.2 documentation.
\url{https://scikit-learn.org/1.7/modules/generated/sklearn.linear_model.ElasticNet.html}.
\end{thebibliography}

\clearpage
\appendix
\section{Feature Definitions and Implementation Details}
\label{app:basic}
Write $c_t,o_t$ for adjusted daily close and open. Raw daily close, open,
high, low, volume, and VWAP are $C_t,O_t,H_t,L_t,V_t,W_t$, respectively;
$v_t$ denotes adjusted volume.
Let $r_t=c_t/c_{t-1}-1$, $D_t=W_tV_t$, and let $\operatorname{mean}_h$ and
$\operatorname{sd}_h$ denote trailing $h$-session mean and sample standard
deviation. A subscript $t-1$ on a rolling statistic means it is lagged by a
session. The target and portfolio are not predictors in this catalogue.

\subsection{The 16 Basic Daily Inputs in M2}
\begin{center}\small
\inlinecaption{table}{All M2 predictors. Every lookback ends no later than signal close.}
\begin{tabular}{p{0.37\linewidth}p{0.56\linewidth}}
\toprule Predictor & Definition\\\midrule
\code{return_1d} & $r_t$\\
\code{return_5d} & $c_t/c_{t-5}-1$\\
\code{return_20d} & $c_t/c_{t-20}-1$\\
\code{momentum_63_skip5} & $c_{t-5}/c_{t-63}-1$\\
\code{momentum_252_skip21} & $c_{t-21}/c_{t-252}-1$\\
\code{day_return} & $c_t/o_t-1$\\
\code{overnight_return} & $o_t/c_{t-1}-1$\\
\code{volatility_5d} & $\operatorname{sd}_5(r)_t$\\
\code{volatility_20d} & $\operatorname{sd}_{20}(r)_t$\\
\code{daily_range} & $(H_t-L_t)/O_t$\\
\code{close_location} & $(2C_t-H_t-L_t)/(H_t-L_t)$\\
\code{ma_deviation_5d} & $c_t/\operatorname{mean}_5(c)_t-1$\\
\code{ma_deviation_20d} & $c_t/\operatorname{mean}_{20}(c)_t-1$\\
\code{relative_volume_5d} & $v_t/\operatorname{mean}_5(v)_{t-1}$\\
\code{relative_volume_20d} & $v_t/\operatorname{mean}_{20}(v)_{t-1}$\\
\code{log_dollar_volume_20d} & $\log\{\operatorname{mean}_{20}(CV)_t\}$\\
\bottomrule
\end{tabular}\end{center}

\clearpage
\subsection{The 28 Daily Inputs Shared by M3 and M4}
Define daytime return $d_t=c_t/o_t-1$, overnight return $n_t=o_t/c_{t-1}-1$,
and market return $r_t^m$ from adjusted SPY closes. The residual is
$e_t=r_t-\beta_{t-1}r_t^m$, where beta uses trailing 126-session
covariance/variance with at least 100 observations, then is lagged one
session. Let $\sigma_t=\operatorname{sd}_{20}(r)_t$.
$\operatorname{rank}_{cs}$ means a percentile rank among eligible stocks.
The six formulaic factors are masked by that same eligibility definition.

{\small
\begin{center}
\inlinecaption{table}{Exact daily-only subset of M4. These 28 columns form M3.}
\begin{tabular}{p{0.40\linewidth}p{0.53\linewidth}}
\toprule Predictor & Definition or construction\\\midrule
\code{reversal_1d} & $-r_t$\\
\code{reversal_5d} & $1-c_t/c_{t-5}$\\
\code{reversal_20d} & $1-c_t/c_{t-20}$\\
\code{momentum_63_skip5} & $c_{t-5}/c_{t-63}-1$\\
\code{momentum_252_skip21} & $c_{t-21}/c_{t-252}-1$\\
\code{day_return_1} & $d_t$\\
\code{day_return_5} & $\exp\{\sum_{j=0}^{4}\log(1+d_{t-j})\}-1$\\
\code{day_return_20} & Same compounding over 20 sessions\\
\code{night_return_1} & $n_t$\\
\code{night_return_5} & $\exp\{\sum_{j=0}^{4}\log(1+n_{t-j})\}-1$\\
\code{night_return_20} & Same compounding over 20 sessions\\
\code{residual_reversal_1} & $-e_t/\sigma_{t-1}$\\
\code{residual_reversal_5} & $-\sum_{j=0}^{4}e_{t-j}/(\sqrt{5}\sigma_{t-1})$\\
\code{volatility_20} & $\sigma_t$\\
\code{illiquidity_20} & $10^6\operatorname{mean}_{20}(|r|/D)_t$\\
\code{relative_dollar_volume} & $D_t/\operatorname{mean}_{20}(D)_{t-1}$\\
\code{normalized_alpha041} & $(\sqrt{H_tL_t}-W_t)/C_t$\\
\code{normalized_vwap_reversal} & $(W_t-C_t)/C_t$\\
\code{close_location_reversal} & $-(2C_t-H_t-L_t)/(H_t-L_t)$\\
\code{alpha002} & Negative 6-session correlation of cross-sectional ranks of
$\Delta_2\log v$ and $(C-O)/O$\\
\code{alpha003} & Negative 10-session correlation of cross-sectional ranks of $O$ and $V$\\
\code{alpha006} & Negative 10-session correlation of $o$ and $v$\\
\code{alpha012} & $-\operatorname{sign}(\Delta v_t)\Delta c_t/a_t$,
where $a_t$ is the adjustment factor\\
\code{alpha041} & $\sqrt{H_tL_t}-W_t$\\
\code{alpha101} & $(C_t-O_t)/(H_t-L_t+0.001)$\\
\code{residual_reversal_x_illiquidity} & Residual reversal (1 session)
times the cross-sectional percentile rank of 20-session illiquidity\\
\code{residual_reversal_x_volume} & Residual reversal (1 session)
times relative dollar volume clipped to $[0.2,5]$\\
\code{gap_reversal_after_disagreement} & $-n_t$ if daytime and overnight
returns have different signs; otherwise zero\\
\bottomrule
\end{tabular}\end{center}}

\clearpage
\subsection{The 18 Intraday Aggregates and Two Tail Interactions}
For regular-session bars $b$, define log return
$u_b=\log(C_b/C_{b-1})$, using the first bar's own open in place of the absent
previous close. Let $U=\sum_b V_b$, $RV=(\sum_b u_b^2)^{1/2}$, $O_s$ be the
session opening bar's open, and $C_s$ the final bar's close. Head and tail
windows are measured from session boundaries, including half days.

{\small
\begin{center}
\inlinecaption{table}{Intraday predictors in M1 and M4. Executable names have an
\texttt{intraday\_} prefix.}
\begin{tabular}{p{0.40\linewidth}p{0.53\linewidth}}
\toprule Name after prefix & Definition\\\midrule
\code{head30} & First-30-minute endpoint close divided by $O_s$, minus one\\
\code{head60} & First-60-minute endpoint close divided by $O_s$, minus one\\
\code{tail30} & $C_s$ divided by the opening price of the last-30-minute window, minus one\\
\code{tail60} & Same for the last 60 minutes\\
\code{head_volume_share} & First-30-minute volume divided by $U$\\
\code{tail_volume_share} & Last-30-minute volume divided by $U$\\
\code{rv} & $RV$\\
\code{downside_share} & $\sum_b u_b^2{\bf1}(u_b<0)/\sum_bu_b^2$\\
\code{skew} & $\sum_bu_b^3/(\sum_bu_b^2)^{3/2}$; a path-shape statistic\\
\code{efficiency} & $|\log(C_s/O_s)|/\sum_b|u_b|$\\
\code{vwap_reversal} & $1-C_s/\{\sum_b VWAP_b V_b/U\}$\\
\code{range} & $(\max_b H_b-\min_b L_b)/O_s$\\
\code{close_location_reversal} & $-(2C_s-\max H-\min L)/(\max H-\min L)$\\
\code{weighted_bar_return} & $\sum_b u_bV_b/U$\\
\code{signed_bar_volume} & $\sum_b\operatorname{sign}(u_b)V_b/U$\\
\code{largest_bar_share} & $\max_b|u_b|/RV$\\
\code{return_volume_corr} & Within-session correlation of $u_b$ and $\log V_b$\\
\code{rth_return} & $C_s/O_s-1$\\
\bottomrule
\end{tabular}\end{center}}

M4 additionally includes \code{tail_reversal_risk_scaled}, equal to the
negative tail-30 return divided by $RV$, and
\code{tail_reversal_x_activity}, equal to negative tail-30 return times the
tail-volume share divided by its lagged 20-session mean (minimum 15
observations). Thus the nested comparison adds 18 aggregates plus two
interactions. Non-finite values remain missing and use the fold's imputer.

\section{Team Contributions}
\begin{table}[htbp]\centering\small
\begin{tabular}{p{0.24\linewidth}p{0.66\linewidth}}\toprule
Member & Contributions\\\midrule
Wan Bicheng & Report writing, code implementation, and video production.\\
Guo Jinghang & Report writing and Model 4 development.\\
Zhao Liyan & Report writing and Models 1 and 2 development.\\
Chen Yifan & Report writing and Model 3 development.\\\bottomrule
\end{tabular}
\end{table}
\end{document}
